\section{Background and Related Work}
\label{sec:background}

\subsection{Propositional Satisfiability and CDCL}
Propositional Satisfiability (SAT) is the problem of determining if there exists an assignment of boolean values to variables that satisfies a given formula. The Boolean satisfiability problem was the first known NP-complete problem \citep{cook1971complexity}. Most modern SAT solvers require the input to be in Conjunctive Normal Form (CNF), meaning the formula is a conjunction of clauses, where each clause is a disjunction of literals. Modern exact SAT solving is dominated by Conflict-Driven Clause Learning (CDCL) algorithms \citep{marques2009conflict}, which extend the classic DPLL backtracking search by inferring and learning new conflict clauses to prune the search space dynamically. Solvers such as Glucose heavily optimize this process by aggressively managing the learned clause database based on Literal Block Distance (LBD) \citep{audemard2009predicting}.

\subsection{At-Most-One (AMO) Encodings}
Translating cardinality constraints---such as ``at most one variable among a set can be true''---into CNF is a critical aspect of SAT modeling. The naive \textbf{Pairwise} (or Binomial) encoding operates without auxiliary variables by generating an exclusion clause for every possible pair, requiring $O(m^2)$ clauses for a group of $m$ variables \citep{frisch2005sat}. 
To mitigate the quadratic clause explosion, auxiliary variables are introduced. The \textbf{Binary} (or Bitwise) encoding maps the $m$ variables to $\lceil \log_2 m \rceil$ auxiliary bits, generating $O(m \log m)$ clauses \citep{frisch2005sat}. The \textbf{Sequential Counter} encoding by \cite{sinz2005towards} mimics a unary adder, bounding both auxiliary variables and clauses to $O(m)$. The \textbf{Commander} encoding \citep{klieber2007solving} recursively partitions the original variables into small groups, electing a ``commander'' variable for each group and generating $O(m)$ auxiliary variables and clauses. Finally, the \textbf{2-Product} encoding structures variables onto a 2D grid, translating the 1D AMO constraint into row and column constraints using $O(\sqrt{m})$ auxiliary variables \citep{chen2010new}. 

\subsection{Constraint Programming and Mixed Integer Programming}
Constraint Programming (CP) provides a higher-level abstraction than SAT by operating directly on integer domains \citep{rossi2006handbook}. Central to CP is the \texttt{AllDifferent} global constraint, which ensures that all variables in a given set take distinct values. Efficient domain filtering algorithms, notably those based on bipartite matching \citep{regin1994filtering}, actively prune the integer domains during search without needing to expand the logic into millions of binary clauses. Tools like OR-Tools CP-SAT \citep{perron2019or} and IBM CP Optimizer leverage these techniques heavily.
Mixed Integer Programming (MIP) relies on linear inequalities over continuous and integer variables \citep{wolsey2020integer}. While powerful for optimization problems, applying MIP solvers like Gurobi or CPLEX to pure feasibility problems like N-Queens translates logical implications into algebraic constraints bounded by linear programming relaxations.
